% year: תשע"ח
% type: מבחן
% semester: א
% moed: ג
% number: 3ב
% points: 10
% categories: mvt, proofs
% source: exams/מבחנים/תשעח/מועד מיוחד תשעח.pdf

\begin{question}
תהי $f$ פונקציה גזירה לכל מספר ממשי $x$, ונניח שיש בדיוק $k$ פתרונות למשוואה $f'(x)=0$. הוכיחו שיש לכל היותר $k+1$ פתרונות למשוואה $f(x)=0$.
\end{question}

\begin{hint}
בין כל שני שורשים של $f$ יש שורש של $f'$ (משפט רול). הוכיחו בדרך השלילה.
\end{hint}

\begin{solution}
נניח בשלילה שלמשוואה $f(x)=0$ יש לפחות $k+2$ פתרונות שונים, ונסמן $k+2$ מהם $x_1<x_2<\dots<x_{k+2}$.

לכל $i=1,\dots,k+1$, הפונקציה $f$ רציפה בקטע $[x_i,x_{i+1}]$ וגזירה בקטע $(x_i,x_{i+1})$ (כי היא גזירה בכל $\R$, ולכן גם רציפה), ו-$f(x_i)=f(x_{i+1})=0$. לפי \textbf{משפט רול} קיימת נקודה $c_i\in(x_i,x_{i+1})$ כך ש-$f'(c_i)=0$.

הקטעים $(x_1,x_2),(x_2,x_3),\dots,(x_{k+1},x_{k+2})$ זרים בזוגות, ולכן $c_1<c_2<\dots<c_{k+1}$ הן $k+1$ נקודות שונות שבהן $f'=0$. זו סתירה להנחה שלמשוואה $f'(x)=0$ יש בדיוק $k$ פתרונות.

לכן למשוואה $f(x)=0$ יש לכל היותר $k+1$ פתרונות. $\blacksquare$
\end{solution}
